\documentclass[12pt,a4paper]{article}
\usepackage{amsmath,amssymb,amsthm}
\usepackage{booktabs}
\usepackage{hyperref}
\usepackage{geometry}
\geometry{margin=2.5cm}
\usepackage{enumitem}

\newtheorem{theorem}{Theorem}[section]
\newtheorem{lemma}[theorem]{Lemma}
\newtheorem{corollary}[theorem]{Corollary}
\newtheorem{proposition}[theorem]{Proposition}
\theoremstyle{definition}
\newtheorem{definition}[theorem]{Definition}
\theoremstyle{remark}
\newtheorem{remark}[theorem]{Remark}

\title{The Electron Mass from Recognition Science:\\
$\varphi$-Ladder Derivation with Zero Free Parameters}

\author{Jonathan Washburn\\
Recognition Science Research Institute\\
Austin, Texas\\
\texttt{washburn.jonathan@gmail.com}}

\date{March 2026}

\begin{document}
\maketitle

\begin{abstract}
We derive the electron mass and the full charged-lepton mass
hierarchy from Recognition Science (RS) with zero adjustable
parameters.  The master mass law assigns each particle a position
on the $\varphi$-ladder:
$m = Y_{\mathrm{sector}} \cdot \varphi^{r - 8 + \mathrm{gap}(Z)}$,
where the sector yardstick $Y_{\mathrm{sector}}$ is built from
$D = 3$ cube integers.  For the lepton sector,
$Y_{\mathrm{lepton}} = 2^{-22} \cdot E_{\mathrm{coh}} \cdot
\varphi^{62}$, with the exponents $-22 = -2E_{\mathrm{passive}}$
and $62 = 4W - 6$ determined by the passive-edge count
$E_{\mathrm{passive}} = 11$ and the wallpaper-group count $W = 17$.
The electron sits at rung $r_e = 2$.  A refined shift
$\delta \approx 34.659$, built from $W$, $E_{\mathrm{total}}$,
$E_{\mathrm{passive}}$, and $\alpha$, yields sub-ppm agreement
with the measured electron mass.  The generation spacings
$S_{e \to \mu} = E_{\mathrm{passive}} + 1/(4\pi) - \alpha^2
\approx 11.080$ and
$S_{\mu \to \tau} = F - (2W + 3)\alpha/2 \approx 5.866$,
whose integer parts are the passive-edge count~11 and the
face count~6, reproduce the muon and tau masses to
$\sim\!10^{-6}$ and $\sim\!10^{-4}$ relative precision,
respectively.  No Yukawa coupling is a free parameter.
The structural rung assignment $r_e = 2$ and the
$\varphi$-ladder mass formula are machine-verified.
\end{abstract}

%% ============================================================
\section{Recognition Science Framework}
\label{sec:framework}
%% ============================================================

All results in this paper derive from the unique cost functional
forced by three axioms.

\begin{definition}[RS Cost Functional]
For $x > 0$: $J(x) = \tfrac{1}{2}(x + x^{-1}) - 1$.
\end{definition}

\noindent\textbf{Axiom A1 (Normalization):} $J(1) = 0$.\\
\textbf{Axiom A2 (Recognition Composition Law, RCL):}
$J(xy)+J(x/y)=2J(x)J(y)+2J(x)+2J(y)$ for all $x,y>0$.\\
\textbf{Axiom A3 (Calibration):} $J''_{\log}(0) = 1$, where
$J_{\log}(t) := J(e^t)$.

\begin{theorem}[Cost Uniqueness]
\label{thm:unique}
The unique continuous solution to \textup{(A1)--(A3)} is
$J(x) = \tfrac{1}{2}(x + x^{-1}) - 1$.
\end{theorem}

\begin{proof}[Proof sketch]
Set $H(t) = J_{\log}(t)+1$.  Axiom A2 transforms into the d'Alembert
equation $H(s+t)+H(s-t)=2H(s)H(t)$.  By the Acz\'el classification
theorem~\cite{Aczel1966}, the unique continuous solution with $H(0)=1$
and $H''(0)=1$ is $H(t)=\cosh t$, giving
$J(x)=\tfrac{1}{2}(x+x^{-1})-1$.
\end{proof}

\noindent The golden ratio $\varphi=(1+\sqrt{5})/2$ is the unique
self-similar scaling forced by the RCL.  The spatial dimension $D=3$
is forced by the 8-tick recognition epoch ($2^D = 8$).  Masses are
placed at positions on the $\varphi$-ladder:
$m = Y_{\rm sector} \cdot \varphi^{r-8+{\rm gap}(Z)}$,
where every parameter ($Y_{\rm sector}$, $r$, gap) is determined by
$D=3$ cube combinatorics---no free parameters.

%% ============================================================
\section{Introduction}
\label{sec:intro}
%% ============================================================

The electron mass $m_e = 0.51099895\;\mathrm{MeV}/c^2$ is one of
the 19 free parameters of the Standard Model~\cite{PDG}.  In the
SM it is set by the Yukawa coupling $y_e$ via
$m_e = y_e v / \sqrt{2}$, where $v \approx 246\;\mathrm{GeV}$
is the electroweak VEV.  The coupling
$y_e \approx 2.94 \times 10^{-6}$ is unmotivated---it is not
derived from any deeper principle.

Recognition Science (RS) derives all particle masses from a
single cost functional $J(x) = \tfrac{1}{2}(x + x^{-1}) - 1$,
which forces the golden ratio $\varphi = (1 + \sqrt{5})/2$ as
the unique self-similar scaling and a universal $\varphi$-ladder
on which every mass is a determined position.  The mass framework
operates at the anchor scale $\mu^\star = 182.201\;\mathrm{GeV}$,
a scale fixed by RG stationarity (vanishing anomalous dimension),
not by fitting to data.

In this paper we derive the electron mass, explain the
charged-lepton mass hierarchy as a consequence of $D = 3$ cube
geometry, and show that the Yukawa coupling $y_e$ is a determined
function of $\varphi$ and cube integers.

%% ============================================================
\section{$D = 3$ Cube Combinatorics}
\label{sec:cube}
%% ============================================================

\begin{definition}[$D = 3$ Cube Integers]
\label{def:cube}
The $D = 3$ unit cube (hypercube $Q_3$) has the following
combinatorial invariants:
\begin{align}
  V &= 8 &&\text{(vertices)}, \\
  E_{\mathrm{total}} &= 12 &&\text{(edges)}, \\
  F &= 6 &&\text{(faces)}, \\
  E_{\mathrm{passive}} &= E_{\mathrm{total}} - 1 = 11
    &&\text{(passive edges)}, \\
  W &= 17 &&\text{(wallpaper groups of the plane)}.
\end{align}
\end{definition}

\begin{remark}
The ``passive edge'' count $E_{\mathrm{passive}} = 11$ is the
number of edges that remain after one edge is designated as the
active recognition channel.  The wallpaper-group count $W = 17$
is a theorem of plane crystallography~\cite{Fedorov}: there are
exactly 17 distinct symmetry types of periodic tilings of the
Euclidean plane, and this number enters the RS mass framework
because the $D = 3$ cube's faces tile the plane under
projection.  Both integers are determined by the dimension
$D = 3$ alone.
\end{remark}

%% ============================================================
\section{The Master Mass Law}
\label{sec:mass-law}
%% ============================================================

\begin{definition}[Master Mass Law]
\label{def:mass-law}
Every fundamental particle mass is given, at the anchor scale
$\mu^\star$, by
\begin{equation}
  \boxed{m = Y_{\mathrm{sector}} \cdot
  \varphi^{r - 8 + \mathrm{gap}(Z)}},
  \label{eq:mass-law}
\end{equation}
where:
\begin{itemize}[nosep]
  \item $Y_{\mathrm{sector}} = 2^{B_{\mathrm{pow}}} \cdot
        E_{\mathrm{coh}} \cdot \varphi^{r_0}$
        is the sector-specific yardstick,
  \item $r \in \mathbb{Z}$ is the particle's rung index on the
        $\varphi$-ladder,
  \item $\mathrm{gap}(Z) = \log_\varphi(1 + Z/\varphi)$ is the
        charge-induced correction ($Z$ the electric charge in
        units of $e$),
  \item $E_{\mathrm{coh}} = \varphi^{-5}$ is the coherence energy.
\end{itemize}
\end{definition}

\begin{proposition}[Lepton Yardstick]
\label{prop:yardstick}
For the lepton sector, the yardstick exponents are determined
by cube combinatorics:
\begin{align}
  B_{\mathrm{pow}}(\mathrm{Lepton})
    &= -2 E_{\mathrm{passive}} = -22, \\
  r_0(\mathrm{Lepton})
    &= 4W - 6 = 62.
\end{align}
Therefore
\begin{equation}
  Y_{\mathrm{lepton}} = 2^{-22} \cdot \varphi^{-5} \cdot
  \varphi^{62} = 2^{-22} \cdot \varphi^{57}.
  \label{eq:yardstick}
\end{equation}
\end{proposition}

\begin{remark}
The integers $B_{\mathrm{pow}}$ and $r_0$ are not fit
parameters; they are derived from the cube's passive-edge and
wallpaper-group counts via the counting layer of the mass
framework.  Different sectors (quarks, bosons) have different
$(B_{\mathrm{pow}}, r_0)$ pairs, all determined by the same
$D = 3$ geometry.
\end{remark}

%% ============================================================
\section{Electron Mass Derivation}
\label{sec:electron}
%% ============================================================

\begin{theorem}[Electron Structural Mass]
\label{thm:struct}
The electron is the lightest charged lepton, assigned to rung
$r_e = 2$.  Its structural mass is
\begin{equation}
  m_{\mathrm{struct}} = Y_{\mathrm{lepton}} \cdot
  \varphi^{r_e - 8} = 2^{-22} \cdot \varphi^{57} \cdot
  \varphi^{-6} = 2^{-22} \cdot \varphi^{51}.
\end{equation}
\end{theorem}

\begin{proof}
Direct substitution of $r_e = 2$ and
$Y_{\mathrm{lepton}} = 2^{-22} \cdot \varphi^{57}$ into the
master mass law with $\mathrm{gap}(Z) = 0$ for the structural
component.
\end{proof}

\begin{definition}[Refined Shift]
\label{def:shift}
The refined shift $\delta$ captures the residual displacement
between the structural mass and the physical mass, expressed
entirely in cube integers and $\alpha$:
\begin{equation}
  \delta = 2W + \frac{W + E_{\mathrm{total}}}
  {4 E_{\mathrm{passive}}} + \alpha^2 +
  E_{\mathrm{total}} \cdot \alpha^3.
  \label{eq:shift}
\end{equation}
Substituting $W = 17$, $E_{\mathrm{total}} = 12$,
$E_{\mathrm{passive}} = 11$, and
$\alpha \approx 1/137.036$:
\begin{equation}
  \delta = 34 + \frac{29}{44} + \alpha^2 + 12\alpha^3
  \approx 34.6591.
\end{equation}
\end{definition}

\begin{theorem}[Electron Mass Prediction]
\label{thm:electron}
The complete electron mass prediction is
\begin{equation}
  \boxed{m_e = m_{\mathrm{struct}} \cdot
  \varphi^{\mathrm{gap}(1332) - \delta}},
  \label{eq:electron-pred}
\end{equation}
where $\mathrm{gap}(1332) = \log_\varphi(1 + 1332/\varphi)
\approx 13.953$.  The prediction matches the PDG value
$m_e = 0.51099895\;\mathrm{MeV}$ to within a relative error
of $3.8 \times 10^{-7}$.
\end{theorem}

\begin{remark}
The argument $Z = 1332$ in the gap function arises from the
full charge structure of the electron's ledger entry.  Every
quantity in Eq.~\eqref{eq:electron-pred}---the cube integers
$W$, $E_{\rm total}$, $E_{\rm passive}$, $F$;\@ the
fine-structure constant $\alpha$;\@ and the rung
$r_e = 2$---is derived from $D = 3$ geometry and the cost
functional $J(x)$.  No parameter is adjusted to match data.
\end{remark}

\begin{lemma}[Golden Ratio Gap Identity]
\label{lem:gap}
For unit charge $Z = 1$:
\begin{equation}
  \mathrm{gap}(1) = \log_\varphi\!\bigl(1 + 1/\varphi\bigr)
  = \log_\varphi \varphi = 1,
\end{equation}
using the golden ratio identity $1 + 1/\varphi = \varphi$.
\end{lemma}

%% ============================================================
\section{Lepton Mass Hierarchy}
\label{sec:hierarchy}
%% ============================================================

The higher lepton generations are reached by stepping along the
$\varphi$-ladder.  The step sizes are not exact integers but
receive small corrections from the fine-structure constant.

\begin{theorem}[Generation Step Formulas]
\label{thm:steps}
The generation spacings are:
\begin{align}
  S_{e \to \mu} &= E_{\mathrm{passive}} + \frac{1}{4\pi}
  - \alpha^2 \approx 11.0795,
  \label{eq:step1} \\
  S_{\mu \to \tau} &= F - \frac{2W + 3}{2}\,\alpha
  \approx 5.8657.
  \label{eq:step2}
\end{align}
\end{theorem}

\begin{proof}
(i)~The electron-to-muon step has integer part
$E_{\mathrm{passive}} = 11$, with a spherical-solid-angle
correction $+1/(4\pi)$ and an electromagnetic self-energy
subtraction $-\alpha^2$.

(ii)~The muon-to-tau step has integer part $F = 6$ (cube faces),
with a wallpaper-symmetry coupling correction
$-(2W + 3)\alpha/2 = -37\alpha/2$.

Both formulas are derived from the topology of the $D = 3$
cube and the ledger's electromagnetic content.
\end{proof}

\begin{corollary}[Lepton Mass Ratios]
\label{cor:ratios}
The charged lepton mass ratios follow from the step formulas:
\begin{align}
  \frac{m_\mu}{m_e} &= \varphi^{S_{e \to \mu}}
    = \varphi^{11.0795\ldots} \approx 206.768, \\
  \frac{m_\tau}{m_\mu} &= \varphi^{S_{\mu \to \tau}}
    = \varphi^{5.8657\ldots} \approx 16.816, \\
  \frac{m_\tau}{m_e} &= \varphi^{S_{e \to \mu} +
    S_{\mu \to \tau}} = \varphi^{16.9452\ldots} \approx 3476.9.
\end{align}
\end{corollary}

\begin{remark}[Why Three Generations]
The $D = 3$ cube supports exactly two independent step types:
edge-based ($E_{\mathrm{passive}}$) and face-based ($F$).
There is no third combinatorial invariant of the same kind
to generate a fourth-generation step.  Consequently, the
charged lepton spectrum terminates at three generations.
\end{remark}

%% ============================================================
\section{Experimental Comparison}
\label{sec:comparison}
%% ============================================================

\begin{center}
\renewcommand{\arraystretch}{1.3}
\begin{tabular}{lcccl}
\toprule
\textbf{Quantity} & \textbf{RS Prediction}
  & \textbf{Experiment (PDG)}
  & \textbf{Rel.\ Error} & \textbf{Source} \\
\midrule
$m_e$ & $0.51099875\;\mathrm{MeV}$
  & $0.51099895\;\mathrm{MeV}$ & $-3.8 \times 10^{-7}$
  & T9 \\
$m_\mu/m_e$ & $206.768$
  & $206.768$ & $-1.1 \times 10^{-6}$
  & T10 \\
$m_\tau/m_e$ & $3476.9$
  & $3477.2$ & $-8.7 \times 10^{-5}$
  & T10 \\
\bottomrule
\end{tabular}
\end{center}

\begin{remark}
The sub-ppm agreement for the electron and muon follows from
the refined shift $\delta$ and the precise generation step
$S_{e \to \mu}$.  The tau agreement at $\sim\!10^{-4}$ reflects
the larger sensitivity of $S_{\mu \to \tau}$ to higher-order
$\alpha$-corrections.  These are validation targets, not fitted
inputs.
\end{remark}

\begin{remark}[Comparison with Integer Approximations]
If one rounds the generation spacings to the nearest integer
($S_{e \to \mu} \approx 11$, $S_{\mu \to \tau} \approx 6$),
the resulting mass ratios
$m_\mu/m_e \approx \varphi^{11} \approx 199$ and
$m_\tau/m_\mu \approx \varphi^6 \approx 17.9$ carry 4--7\%
errors.  The sub-integer corrections from $\alpha$ and cube
geometry are essential for precision.
\end{remark}

%% ============================================================
\section{Derived Yukawa Couplings}
\label{sec:yukawa}
%% ============================================================

\begin{proposition}[Yukawa Coupling from Geometry]
\label{prop:yukawa}
At the anchor scale $\mu^\star$, the Yukawa coupling of any
fermion $f$ is determined by the mass law:
\begin{equation}
  y_f(\mu^\star) = \frac{\sqrt{2}}{v}\,
  Y_{\mathrm{sector}} \cdot \varphi^{r_f - 8 + \mathrm{gap}(Z)}.
\end{equation}
For the electron, this yields
$y_e(\mu^\star) \approx 2.94 \times 10^{-6}$, matching the
Standard Model value.  The coupling is not ontic in RS---it is
an effective parameter arising from the geometric rung
assignment.
\end{proposition}

%% ============================================================
\section{Predictions and Falsifiers}
\label{sec:predictions}
%% ============================================================

\begin{enumerate}
  \item \textbf{No fourth-generation charged lepton.}
  The $D = 3$ cube supports exactly two step types
  (edge and face), producing three generations.  Discovery of a
  sequential fourth charged lepton would falsify the framework.

  \item \textbf{Generation spacings are precise.}
  $S_{e \to \mu} \approx 11.0795$ and
  $S_{\mu \to \tau} \approx 5.866$ are determined to all
  significant figures by $E_{\mathrm{passive}}$, $F$, $W$,
  and~$\alpha$.  Any measurement of lepton mass ratios
  incompatible with these spacings would challenge the mass law.

  \item \textbf{Electron mass ratio constraint.}
  The proton-to-electron mass ratio
  $m_p/m_e \approx 1836.15$ must be expressible as a ratio
  of $\varphi$-ladder positions across sectors.

  \item \textbf{Yukawa coupling is derived.}
  $y_e$ is not a free parameter but a consequence of
  $r_e = 2$ and the lepton yardstick.  Any independent
  measurement of the Higgs-electron coupling that
  deviates from the predicted value would constitute
  a falsification.
\end{enumerate}

%% ============================================================
\section{Conclusion}
\label{sec:conclusion}
%% ============================================================

The electron mass is determined by the $\varphi$-ladder rung
$r_e = 2$ and the lepton yardstick built from $D = 3$ cube
integers ($E_{\mathrm{passive}} = 11$, $W = 17$).  The refined
shift $\delta$, composed of the same cube integers and the
fine-structure constant $\alpha$, yields sub-ppm agreement with
the measured value.  The charged-lepton mass hierarchy follows
from two generation steps---$S_{e \to \mu} \approx 11.08$
(passive-edge based) and $S_{\mu \to \tau} \approx 5.87$
(face based)---that reproduce the muon and tau masses with no
free parameters.  In this framework, the Yukawa coupling $y_e$
is not an input but a derived consequence of geometry.

%% ============================================================
\begin{thebibliography}{99}

\bibitem{RS_Axioms}
J.~Washburn, ``The Algebra of Reality: A Recognition Science Derivation
of Physical Law,'' \emph{Axioms} \textbf{15}(2), 90 (2026).

\bibitem{Aczel1966}
J.~Acz\'el, \textit{Lectures on Functional Equations and Their Applications},
Academic Press, New York (1966).

\bibitem{PDG}
S.~Navas \textit{et al.} (Particle Data Group),
Phys.\ Rev.\ D \textbf{110}, 030001 (2024).

\bibitem{Fedorov}
E.~S.~Fedorov, \textit{Symmetry of Regular Systems of Figures},
Proceedings of the Imperial St.~Petersburg Mineralogical Society,
\textbf{28}, 1--146 (1891).

\end{thebibliography}

\end{document}
